\section{2. Related Work}

\subsection{Traditional Light Field Depth Estimation Methods}

Traditional light field depth estimation methods have emerged as a fundamental paradigm, extensively leveraging the rich angular information to infer scene geometry. Let a 4D light field be parameterized by spatial coordinates $(x, y)$ and angular coordinates $(u, v)$. Wanner et al.  proposed a globally consistent depth labeling method that utilizes the structure tensor on epipolar plane image (EPI)s (EPIs) to extract geometric disparities. Specifically, the orientation of the linear trajectories in EPIs is explicitly formulated to derive depth, followed by a Markov Random Field (MRF) optimization to ensure spatial coherence. This geometric approach demonstrates remarkable accuracy in controlled environments with rich textures. However, it fundamentally relies on the global Lambertian assumption. When encountering glossy or metallic surfaces, the viewpoint-dependent intensity variations dramatically disrupt the linear structures in EPIs, leading to severe aliasing bimodal distributions and completely invalidating the slope calculation.

While EPI-based methods primarily focus on geometric slopes, other approaches extend to photometric cues, such as defocus and correspondence, to mitigate the impact of non-Lambertian reflectances. Tao et al. [1] proposed a clustering method that eliminates specular pixels when enforcing photo-consistency by combining defocus and correspondence cues from light-field cameras. By analyzing the focal stack, their method attempts to weight or discard specular highlights to recover the underlying geometry. Unfortunately, they attempt a binary classification of pixels into either Lambertian or specular, which cannot handle general glossy surfaces with complex spatially-varying BRDFs (SVBRDFs). Moreover, the defocus cues often become ambiguous in textureless regions, making the depth estimation unreliable for mixed materials and arbitrarily failing on general non-Lambertian objects.

To further address complex reflectances beyond simple binary classifications, subsequent studies have incorporated traditional physical reflection models into multi-view stereo paradigms. For instance, methods integrating the dichromatic model [2] attempt to separate diffuse and specular components before performing stereo matching. By formulating the reflectance as a linear combination of body and surface reflections, these methods aim to accurately recover the shape of glossy objects. Nonetheless, the dichromatic model fails to hold for general translucent or scattering materials, where light undergoes complex subsurface scattering. Additionally, estimating physical parameters alongside depth significantly increases the computational complexity, rendering these methods impractical for high-resolution light field data and mostly limiting them to synthetic or highly controlled scenarios.

In short, many traditional depth estimation methods for light-field cameras have been proposed. However, most of them rely on the global Lambertian assumption and work poorly on non-Lambertian surfaces. Specifically, these methods exhibit two critical limitations. First, they severely depend on the global Lambertian assumption; in non-Lambertian regions such as specular or scattering surfaces, the violation of the viewpoint consistency assumption causes the EPI slope calculation to completely fail. Second, traditional physical models exhibit extremely high computational complexity and struggle to adaptively process complex material distributions in real-world mixed scenes (e.g., urban environments), lacking a robust mechanism to uniformly handle Lambertian and non-Lambertian regions. In contrast, our model abandons the global Lambertian assumption and proposes a component-aware multi-strategy adaptive depth estimation and fusion framework. Specifically, we employ EPI for diffuse regions, photometric stereo and neighborhood propagation for specular regions, and scattering model correction for scattering regions. Furthermore, unlike previous methods that rely on black-box classifiers or simple spatial textures, we introduce an MRI-like angular frequency analysis mechanism and a dual-mask modulated physical-driven feature representation, enabling highly interpretable, robust, and unified depth estimation across complex real-world scenes.

\subsection{Deep Learning-based Light Field Depth Estimation and Feature Representation}

Following the limitations of handcrafted geometric and photometric cues, deep learning-based approaches have emerged as a dominant paradigm for light field depth estimation, extensively leveraging convolutional neural networks (CNNs) and Transformers to implicitly learn complex feature representations. Shin et al.  proposed EPINet, a pioneering fully convolutional network that directly extracts depth cues from epipolar plane image (EPI)s (EPIs). By treating the EPI as a 2D image and applying multi-scale CNN architectures, EPINet effectively captures the linear slope variations corresponding to scene disparities. While this EPI-centric approach demonstrates remarkable accuracy on standard Lambertian datasets, it primarily focuses on local 2D EPI patches, often failing to capture the global spatial-angular context. Consequently, when encountering non-Lambertian surfaces where the linear EPI structures are severely distorted by specular highlights, EPINet struggles to distinguish between geometric disparities and reflectance-induced intensity shifts, leading to erroneous depth predictions.

To address the restricted receptive field of purely EPI-based methods, subsequent research has extended to multi-view and sub-aperture fusion networks. Liang et al. [3] introduced LFNet, which formulates depth estimation by constructing a cost volume from sub-aperture image arrays. Specifically, LFNet extracts spatial features from individual sub-aperture images and aggregates them across different viewpoints to evaluate photo-consistency, significantly improving depth accuracy in textureless regions. Moreover, by incorporating a multi-scale feature aggregation module, it enhances the robustness of disparity prediction. However, this cost volume-based fusion heavily relies on the strict photo-consistency assumption across viewpoints. In the presence of spatially-varying BRDFs (SVBRDFs), the viewpoint-dependent intensity variations dramatically violate this assumption, leading to erroneous cost volume matching and severe depth artifacts on glossy or metallic objects.

In contrast to relying solely on EPI slopes or sub-aperture photo-consistency, more recent architectures have attempted to explicitly integrate defocus cues to compensate for the failure of correspondence matching in non-Lambertian regions. For instance, recent dual-branch networks [4] propose a defocus-EPI dual-branch framework, where one branch extracts geometric features from EPIs while the other captures defocus blur cues from focal stacks. The features from both branches are subsequently fused via an attention mechanism to refine the final depth map. This multi-cue strategy theoretically provides complementary information, as defocus blur is generally less sensitive to specular reflections than epipolar lines. Unfortunately, in complex illumination conditions, the defocus branch often fails to learn discriminative depth features due to the entanglement of defocus blur and specular highlights. As a result, the additional branch frequently degenerates into an ineffective module, contributing marginally to the final performance and failing to fundamentally resolve the feature representation breakdown in non-Lambertian areas.

In short, while existing deep learning-based methods have significantly advanced light field depth estimation, they share fundamental limitations when applied to non-Lambertian scenes. First, the black-box nature of implicit feature extraction lacks physical interpretability. In non-Lambertian regions, these networks are highly susceptible to degrading into "texture copying," where the model arbitrarily outputs the input high-frequency textures rather than inferring the true underlying depth structure. Second, merely appending auxiliary branches (such as the defocus branch) without physical guidance proves insufficient; these branches often struggle to learn valid depth features under complex lighting, ultimately becoming dead ends that fail to resolve the root cause of feature representation failure. To overcome these bottlenecks, our work diverges from purely data-driven paradigms by proposing a dual-mask modulated physics-driven light field feature representation model. By introducing a medium mask and an angular direction mask, and coupling them with RTF rank analysis, we explicitly transform implicit black-box features into physically meaningful mask modulations. This physics-aware mechanism effectively mitigates the texture copying issue, providing a robust and interpretable foundation for unified depth estimation across Lambertian, non-Lambertian, and complex mixed scenes.

\subsection{Deep Learning-based Material Parsing and Non-Lambertian Depth Estimation}

While the aforementioned cost volume and EPI-based networks implicitly learn feature representations, explicitly parsing material properties is crucial for handling non-Lambertian surfaces. To this end, Wang et al. [5] proposed a CNN-based material classification network that extracts spatial-angular features from light field patches to identify surface reflectance properties. By formulating material recognition as a multi-class classification problem, their model demonstrates remarkable accuracy on synthetic datasets with distinct, uniform material categories. However, this approach primarily relies on local spatial textures within small patches (e.g., $9 \times 9$). Such restricted receptive fields often lack sufficient angular and spatial context to disambiguate complex spatially-varying BRDFs (SVBRDFs), making them unreliable for general glossy surfaces in real-world scenarios.

To explicitly mitigate the interference of specular highlights, subsequent research has explored reflection separation and joint estimation frameworks. Li et al. [6] introduced a deep reflection separation network coupled with a joint depth estimation module. Leveraging the dichromatic reflection model, this method attempts to decompose the input light field into diffuse and specular components, subsequently feeding the recovered diffuse component into a standard cost volume network for depth prediction. While this decomposition strategy significantly alleviates the impact of severe highlights, the separation network is usually trained with pseudo labels or in a weakly supervised manner due to the scarcity of ground-truth reflection components. Consequently, the recovered diffuse images often exhibit in-painting artifacts and aliasing bimodal distributions, which dramatically degrade the downstream depth accuracy.

In contrast to explicit physical decomposition, other approaches attempt to handle non-Lambertian distortions implicitly within the depth estimation pipeline. More recently, Chen et al. [7] developed a specular-aware depth estimation framework specifically tailored for non-Lambertian light fields. Instead of separating reflection components, they designed an attention-based feature fusion module that implicitly down-weights EPI regions distorted by specular highlights, thereby preserving the geometric consistency of the underlying Lambertian structures. Although this implicit weighting properly addresses surfaces with mild glossiness, it fundamentally struggles with highly reflective or metallic regions where the diffuse component is entirely absent. In such cases, the attention mechanism arbitrarily suppresses valid geometric cues, leading to severe depth ambiguities and erroneous predictions.

Despite these advancements, existing deep learning-based material parsing and non-Lambertian depth estimation methods exhibit critical limitations. First, when relying on spatial textures or black-box CNN classifiers (e.g., a standard 3-class CNN) for material parsing, the local patches (e.g., $9 \times 9$) contain insufficient angular and spatial information, resulting in extremely low classification accuracy (e.g., merely 24.6\%) and failing to provide reliable physical priors. Second, existing methods fail to establish an effective physical cascade between material parsing and depth estimation. They lack efficient and interpretable material analysis from the perspectives of the frequency domain and physical optics, and do not achieve adaptive deep fusion based on material components. To address these fundamental bottlenecks, we propose an MRI-inspired angular frequency-domain analysis mechanism for light field material parsing, utilizing 2D-DFT to efficiently and interpretably classify materials in the frequency domain. Coupled with a dual-mask modulated physical-driven feature representation and a component-aware adaptive fusion architecture, our approach achieves high-precision unified processing of Lambertian, non-Lambertian, and complex mixed scenes within a single model, which has not been achieved by previous light field approaches.

In summary, existing light field depth estimation methods, whether relying on implicit feature learning or explicit black-box material classification, either fail under violated view-consistency assumptions in non-Lambertian regions or struggle to disentangle complex mixed materials due to restricted spatial-angular receptive fields. Inspired by these observations, we propose a unified dual-mask physical approach that addresses the aforementioned limitations by integrating an MRI-inspired angular frequency analysis for interpretable material parsing, a dual-mask modulation mechanism for physically-driven feature representation, and a component-aware adaptive fusion framework to robustly estimate depth across diverse Lambertian, non-Lambertian, and mixed real-world scenes.